%*****************************************
\chapter{Conclusions}\label{ch:conclusions}
%*****************************************

% \begin{flushright}
%      \begin{minipage}{0.73\textwidth}
%     \raggedleft
%     % \slshape\small
%     % When you can measure what you are speaking about, and express it in numbers, you know something about it; but when you cannot measure it, when you cannot express it in numbers, your knowledge is of a meagre and unsatisfactory kind.\\ \medskip
%     % {\normalfont --- Lord Kelvin}
%     \slshape\small\color{mid-gray}
%     Essentially, all models are wrong, but some are useful.\\ \medskip
%     {\normalfont ---\, George E. P. Box}
%     \end{minipage}
% \end{flushright}
% \vspace{1.5em}

\begin{flushquote}
    {Essentially, all models are wrong, but some are useful.}
    {George E. P. Box}
\end{flushquote}

Posterior consistency of the Bayes factor in a point-null testing situation was established by \citet{dass2004note}, but their abstract approach cannot provide posterior contraction rates. As a novel contribution, we show that under \(H_0\) the posterior probability of the alternative goes to zero at a polynomial rate, while under \(H_1\) the posterior probability of the null decays exponentially fast.

\section{Future work and open research lines}

\note{Inferencia bayesiana usando GP como prior --> citar luis?}
\note{Estimación puntual bayesiana? selección variables? (similar a chapter 2)}

%*****************************************
%*****************************************
%*****************************************
%*****************************************
%*****************************************